% Section: Method — IA-SASRec
% =====================================================================
% Lifted from ia_sasrec.md §2--4. Three variants, learnable lambda per
% layer (init 1.0), normalisation as a hyperparameter.
% =====================================================================
\section{IA-SASRec: Intensity-Aware Self-Attention}
\label{sec:method}

\subsection{Background: SASRec attention}
\label{sec:method:background}

SASRec~\citep{kang2018sasrec} encodes a user's left-padded history
$s = [i_1, \dots, i_T]$ (sequence length $T$) as embeddings
$E \in \mathbb{R}^{T \times d}$, where $d$ is the embedding dimension.
Each transformer block projects $E$ to queries, keys, and values
$Q, K, V \in \mathbb{R}^{T \times d}$, and computes scaled dot-product
self-attention with a causal mask $M \in \{0, -\infty\}^{T \times T}$:
\begin{equation}
\mathrm{Attention}(Q, K, V) = \mathrm{softmax}\!\left(\tfrac{QK^{\top}}{\sqrt{d}} + M\right) V.
\label{eq:sasrec}
\end{equation}
Only an item's identity (via the embedding lookup) and its position (via the
positional embedding) enter equation~\eqref{eq:sasrec}; the
strength of each interaction (a 5-star rating, hours played,
or dwell time) is not used. IA-SASRec re-introduces that signal.

Let $w \in \mathbb{R}^T$ be the per-position intensity vector for the
user's history (after the normalisation step described below). Define
the intensity matrix $M_W \in \mathbb{R}^{T \times T}$ with
$M_W[q, k] = w_k$, broadcasting intensity along the key axis,
which is the axis the softmax normalises.

\subsection{Three injection points}
\label{sec:method:variants}

The attention expression~\eqref{eq:sasrec} has three algebraic positions
where $M_W$ can be threaded in: as a pre-softmax additive bias, as a
pre-softmax multiplicative scale, or as a post-softmax attention reweighting.
We instantiate all three. Each adds one learnable scalar
$\lambda \in \mathbb{R}$ per attention layer, initialised to $1.0$, so
that $\lambda = 0$ collapses the variant exactly to
SASRec~\eqref{eq:sasrec}. This learnable gate is the central design
discipline of IA-SASRec.

\paragraph{IA-SASRec-Add (additive logit bias).}
Intensity enters as an explicit bias on the raw attention logits:
\begin{equation}
\mathrm{Att}_{\mathrm{Add}} = \mathrm{softmax}\!\left(\tfrac{QK^{\top}}{\sqrt{d}} + \lambda\, M_W + M\right) V.
\label{eq:add}
\end{equation}
A high-intensity key $k$ inflates column $k$ of the logit matrix and is
forced to receive large attention probability, irrespective of semantic
similarity to the query. As $\lambda \to \infty$ the variant approaches
a hard intensity-argmax.

\paragraph{IA-SASRec-Mul (multiplicative logit scaling).}
Instead of an additive bias, the semantic similarity itself is scaled
by intensity:
\begin{equation}
\begin{aligned}
\mathrm{Att}_{\mathrm{Mul}} = \mathrm{softmax}\Big(
  & \tfrac{QK^{\top}}{\sqrt{d}} \odot \big(1 + \lambda(M_W - 1)\big) \\
  & {} + M \Big)\, V,
\end{aligned}
\label{eq:mul}
\end{equation}
with $\odot$ element-wise. The $(w_k - 1)$ parameterisation interpolates
between two endpoints: at $\lambda = 0$ the multiplier is $1$ and the
variant collapses to SASRec; at $\lambda = 1$ the multiplier is exactly
$w_k$, recovering the original hard-gatekeeper formulation. The causal
mask is added after scaling so padded keys remain at $-\infty$.

\paragraph{IA-SASRec-Val (post-softmax attention reweighting).}
The standard softmax computes the attention distribution; intensity then
reweights each key position's attention probability before the value
mix. Equivalently, since the factor depends only on $k$, this is
algebraically identical to pre-scaling each value vector $V_k$ by its
intensity: $\big(A \odot (1{+}\lambda(M_W{-}1))\big)\,V = A\,(D_\lambda V)$,
where $D_\lambda = \mathrm{diag}(1{+}\lambda(w{-}1))$. Hence the name
\emph{Val}: the operation acts at the value-mixing stage of attention.
\begin{equation}
\begin{aligned}
\mathrm{Att}_{\mathrm{Val}} = \bigl[\,
  & \mathrm{softmax}\!\left(\tfrac{QK^{\top}}{\sqrt{d}} + M\right) \\
  & \odot \bigl(1 + \lambda(M_W - 1)\bigr) \bigr]\, V.
\end{aligned}
\label{eq:val}
\end{equation}
Unlike Mul, which acts on raw logits and can be diluted by the
softmax normalisation, Val acts directly on the attention
probabilities that mix the values. We deliberately do not renormalise
the reweighted distribution: any magnitude inflation introduced by
$1{+}\lambda(w_k{-}1) > 1$ is absorbed by the post-attention layer
normalisation~\citep{ba2016layernorm} in the residual block, while the
intensity-driven \emph{direction} change in the output vector --- the
part that actually re-ranks candidates --- passes through unaffected.

Equation~\eqref{eq:val} is a breaking change from a previous formulation
of Val, which multiplied the post-attention context by
query-position intensity. At prediction time only the last query is
read out, so every candidate score was rescaled by the same constant
and the ranking was identical to vanilla SASRec: the variant could
not reorder candidates and is not a meaningful test of intensity
injection.

\subsection{The $\lambda$ fallback}
\label{sec:method:escape}

At $\lambda = 0$ each variant equals SASRec bit-for-bit (verified by
unit tests on the attention math). A fair learner that finds intensity
not informative should therefore drive $\lambda \to 0$ and recover the
baseline. This is the load-bearing prediction we test empirically: if
the fallback holds, learned $\lambda$
should be small wherever IA-SASRec fails to improve over SASRec, and
large only where it helps.

\subsection{Intensity normalisation}
\label{sec:method:norm}

Raw intensities are heterogeneous: Steam hours range from $0.1$ to
$1000+$, while ratings live in $\{1,\dots,5\}$. Feeding raw values
through a softmax causes severe value collapse, so normalisation is
essential rather than cosmetic. We expose four modes as a single
hyperparameter \texttt{intensity\_norm} selected by Optuna~\citep{akiba2019optuna}:
\texttt{log1p\_minmax} (default; per-user $\log(1+w)$ then divide by
the masked max), \texttt{minmax} (per-user with a floor $f = 0.1$ so
the weakest real signal stays distinguishable from padding),
\texttt{zscore}, and \texttt{none} (sanity-check ablation). The floor
$f$ is non-cosmetic: without it the least-intense real item maps to
exactly $0$ and becomes indistinguishable from padding, silently
discarding one interaction per sequence in Mul and Val. After every
mode, padding positions are forcibly zeroed.

\subsection{Parameter cost}
\label{sec:method:cost}

Each variant adds exactly $n_{\mathrm{layers}}$ learnable scalars
total ($2$--$3$ in our configurations), strictly less than any
prior side-information injection. Table~\ref{tab:variants} summarises
the three formulations.

\begin{table*}[t]
\centering
\small
\caption{The three IA-SASRec variants. $\lambda$ is a learnable scalar per layer (init $1.0$); $\lambda = 0$ collapses each variant to SASRec.}
\label{tab:variants}
\begin{tabular}{ll}
\toprule
Variant & Mechanism \\
\midrule
Add & $\mathrm{softmax}\!\left(\dfrac{QK^{\top}}{\sqrt{d}} + \lambda\, M_W + M\right)\! V$ \\[1.2em]
Mul & $\mathrm{softmax}\!\left(\dfrac{QK^{\top}}{\sqrt{d}} \odot \bigl(1 + \lambda(M_W - 1)\bigr) + M\right)\! V$ \\[1.2em]
Val & $\left[\,\mathrm{softmax}\!\left(\dfrac{QK^{\top}}{\sqrt{d}} + M\right) \odot \bigl(1 + \lambda(M_W - 1)\bigr)\right] V$ \\[0.6em]
\bottomrule
\end{tabular}
\end{table*}
